\chapter{From Local Conservation to Einstein's Field Equations}

Einstein's field equations are simple enough to write in one line and
notoriously unpleasant to expand in components. A metric has many Christoffel
symbols; the Christoffel symbols contain derivatives of the metric; and the
curvature tensor contains both derivatives and products of Christoffel
symbols. Even before one tries to solve anything, the calculation can occupy
pages. For a particular metric, symbolic software is often the sensible way
to manage that bookkeeping.

The principles are much cleaner than the computation. We will therefore not
begin by attacking the equations component by component. We will ask what
kind of equation gravity requires. The argument starts with Newton's theory,
passes through local conservation, and ends by discovering the one
combination of curvature tensors that can consistently equal the flow and
density of energy and momentum. Throughout the chapter we use \(c=1\).

\section{The Newtonian Two-Way Street}

Gravity has two complementary instructions. The gravitational field tells
matter how to move, while matter tells the gravitational field what to be.
This two-way street is already visible in Newtonian mechanics.

Let \(\phi(\mathbf{x},t)\) be the Newtonian gravitational potential. A test
body of mass \(m\) feels the force
\begin{equation}
 \mathbf{F}=-m\nabla\phi.
\end{equation}
Combining this with \(\mathbf{F}=m\mathbf{a}\) cancels the test mass:
\begin{equation}
 \mathbf{a}=-\nabla\phi.
\end{equation}
This is the first direction of the street. Once the field is known, it tells
every freely falling body how to accelerate, independently of that body's
mass.

The reverse direction is Poisson's equation,
\begin{equation}
 \nabla^2\phi=4\pi G\rho,
 \qquad
 \nabla^2=
 \frac{\partial^2}{\partial x^2}
 +\frac{\partial^2}{\partial y^2}
 +\frac{\partial^2}{\partial z^2},
 \label{eq:poisson}
\end{equation}
where \(\rho(\mathbf{x},t)\) is mass per unit volume. The factor \(4\pi\)
could be moved into the definition of \(G\), but it would only reappear
elsewhere. Equation~\eqref{eq:poisson} says that the distribution of matter
is the source of the potential.

For a spherically symmetric body of total mass \(M\), the exterior solution is
\begin{equation}
 \phi(r)=-\frac{GM}{r}.
 \label{eq:newton-potential}
\end{equation}
The density may be concentrated near the center, spread uniformly, or tapered
smoothly toward the surface. Spherical symmetry and the total mass are all
that the exterior solution remembers. Outside the matter,
\(\nabla^2\phi=0\); differentiating \(1/r\) then gives the inverse-square
force. If we care only about the exterior, we may shrink the source to a
point and use the same expression everywhere except at that point.

\begin{classroomqa}
\audiencequestion{Must the density be constant inside the spherical body and
then jump discontinuously to zero at its boundary?}
\lecturerresponse{No. It must be spherically symmetric for this argument, but
it need not be constant and it need not jump. It may decrease smoothly to
zero. None of those interior details changes the exterior potential, which
depends only on the total mass.
\lecturetimestamp{00:09:13--00:09:37}}
\end{classroomqa}

This field formulation is more local-looking than a direct force between
distant masses, but Newtonian gravity still contains instantaneous action at
a distance. Move the source, and the Newtonian potential changes everywhere
at once. General relativity must replace that instantaneous response as well
as the force law.

\section{A Weak-Field Clue, Not Yet an Equation}

The relativistic version of the first instruction is already known: once the
spacetime metric is given, freely falling particles follow spacetime
geodesics. We now need the relativistic replacement for Poisson's equation,
the instruction by which matter determines the metric.

The Schwarzschild geometry offers a clue. Its time-time component is
\begin{equation}
 g_{00}=1-\frac{2GM}{r}.
\end{equation}
Using Equation~\eqref{eq:newton-potential}, this becomes
\begin{equation}
 g_{00}=1+2\phi.
\end{equation}
Consequently,
\begin{equation}
 \nabla^2g_{00}=2\nabla^2\phi=8\pi G\rho.
 \label{eq:weak-mnemonic}
\end{equation}

Equation~\eqref{eq:weak-mnemonic} must be taken with a grain of salt. It is a
mnemonic, not Einstein's equation. It uses a preferred component of the
metric and a nonrelativistic source. Still, it displays the texture we are
looking for: a matter density on one side and second derivatives of geometry
on the other. It is the rough bridge from Poisson's equation to a covariant
field equation.

The right-hand side now needs closer attention. In relativistic units mass is
energy, so \(\rho\) should first be understood as an energy density. But an
energy density is not a complete relativistic object. To see what accompanies
it, we begin with a simpler conserved quantity.

\section{Charge Density, Current, and Local Conservation}

Let \(Q\) be electric charge and let \(\sigma\) denote charge density:
\begin{equation}
 \sigma=\lim_{\Delta V\to0}\frac{\Delta Q}{\Delta V}.
\end{equation}
Density asks how much charge occupies a small spatial volume at one instant.
Current asks a different question. Place an infinitesimal oriented window in
space and count how much charge crosses it per unit area per unit time. The
three possible orientations give the spatial current components \(J^i\).

The dimensional comparison is instructive. Charge density is charge divided
by three spatial lengths. Current is charge divided by two spatial lengths
and one time interval. When \(c=1\), time and length have the same units, so
density and current have matching dimensions. They fit together as a
four-vector:
\begin{equation}
 J^\mu=(J^0,J^1,J^2,J^3),
 \qquad J^0=\sigma.
\end{equation}

Charge conservation is not merely a statement that the total charge in the
universe stays fixed. It is local. Charge cannot vanish from the classroom
and instantaneously reappear at Alpha Centauri. If it leaves a region, it
must cross the region's boundary. For a small box, the decrease of charge
inside equals the outward flux through its walls. Gauss's theorem gives
\begin{equation}
 \dot{\sigma}+\nabla\cdot\mathbf{J}=0,
\end{equation}
or, using the summation convention,
\begin{equation}
 \partial_\mu J^\mu=0.
 \label{eq:charge-continuity}
\end{equation}
In curved coordinates the corresponding tensor equation is written with a
covariant derivative:
\begin{equation}
 D_\mu J^\mu=0.
\end{equation}
The purpose of the covariant derivative is not decoration. It makes the
statement meaningful in any coordinate system.

\section{Why Energy and Momentum Need Two Indices}

Charge itself is invariant. Observers in relative motion agree on the charge
of an electron, although they do not agree on the charge density or current.
A moving observer sees Lorentz-contracted volumes and may see a wind of
charges where another observer sees charges at rest. Density and current
therefore mix under a change of frame even though total charge does not.

Energy and momentum are more complicated. Even their totals are
frame-dependent. One observer may see a particle at rest, carrying only its
rest energy and no spatial momentum. Another observer walking past sees both
kinetic energy and momentum. The four quantities form the momentum
four-vector
\begin{equation}
 P^\nu=(P^0,P^1,P^2,P^3),
 \qquad P^0=E.
\end{equation}
Each component \(P^\nu\) is locally conserved. Each therefore has both a
density and a flow. One index identifies the conserved quantity; a second
index identifies density or direction of flow.

Start with energy. Its density is \(T^{00}\). Energy can also cross a window,
so \(T^{0i}\) is the flow of energy in the \(x^i\) direction. Now take the
\(i\)-component of momentum. Its density is \(T^{i0}\), and its flow in the
\(x^j\) direction is \(T^{ij}\). In this convention:
\begin{itemize}
 \item \(T^{00}\) is energy density;
 \item \(T^{0i}\) is energy flux in direction \(i\);
 \item \(T^{i0}\) is density of the \(i\)-component of momentum;
 \item \(T^{ij}\) is flux of \(i\)-momentum through a surface normal to \(j\).
\end{itemize}
The spatial block \(T^{ij}\) will later be recognized as stress: its diagonal
entries include pressures, while its off-diagonal entries describe shear.

\begin{classroomqa}
\audiencequestion{When \(T^{m0}(x)\) is called the density of the
\(m\)-component of momentum, does \(x\) itself mean the \(m\)-component of
position?}
\lecturerresponse{No. The argument \(x\) denotes the spacetime position at
which the density is evaluated; it stands for all the coordinates. The index
\(m\) labels which momentum component is being counted.
\lecturetimestamp{00:41:07--00:41:27}}
\end{classroomqa}

For each fixed \(\nu\), local conservation gives a continuity equation:
\begin{equation}
 D_\mu T^{\nu\mu}=0.
 \label{eq:stress-continuity}
\end{equation}
The four values of \(\nu\) give one equation for energy and three for
momentum. Together the densities and flows form the energy-momentum tensor.

\begin{figure}[htbp]
\centering
\includegraphics[width=0.84\textwidth]{lecture_09_figure_02.png}
\caption{The developing energy-momentum array on the board. The first index
selects energy or a momentum component; the second selects density or flow.}
\lectureframe{00:44:30}
\end{figure}

The full array is
\begin{equation}
 T^{\mu\nu}=
 \begin{pmatrix}
 T^{00}&T^{01}&T^{02}&T^{03}\\
 T^{10}&T^{11}&T^{12}&T^{13}\\
 T^{20}&T^{21}&T^{22}&T^{23}\\
 T^{30}&T^{31}&T^{32}&T^{33}
 \end{pmatrix}.
\end{equation}
A further theorem of relativistic mechanics and field theory makes the
stress-energy tensor relevant to gravity symmetric:
\begin{equation}
 T^{\mu\nu}=T^{\nu\mu}.
 \label{eq:stress-symmetry}
\end{equation}
In particular, energy flux and momentum density are the corresponding
off-diagonal entries viewed from opposite sides. Symmetry also lets us write
Equation~\eqref{eq:stress-continuity} in the familiar form
\begin{equation}
 D_\mu T^{\mu\nu}=0.
 \label{eq:stress-conservation}
\end{equation}

\begin{classroomqa}
\audiencequestion{Why did energy-momentum acquire an extra index when charge
current needed only one?}
\lecturerresponse{The current index already distinguishes density from flow
direction. Charge needs no additional label because the conserved charge is
a scalar. Energy and the three momenta are themselves four components, so
the first index says which conserved quantity is involved and the second says
whether it is a density or a flow.
\lecturetimestamp{00:45:31--00:47:24}}
\end{classroomqa}

\begin{classroomqa}
\audiencequestion{Is the entire \(T^{\mu\nu}\) matrix symmetric, not merely
the energy-flux and momentum-density entries?}
\lecturerresponse{Yes. The entire matrix is symmetric. The detailed meaning
of its spatial entries, including pressure and stress, will be developed
later.
\lecturetimestamp{00:47:32--00:47:54}}
\end{classroomqa}

Energy density was therefore only one corner of a larger object. The source
of gravity cannot consistently be \(T^{00}\) alone; it must involve the full
tensor.

\section{A Law Cannot Select One Component}

The need for a tensor equation can be seen without any curvature calculation.
Suppose \(A_i\) and \(B_i\) are vectors and a proposed law says only
\(A_3=B_3\). If that equality is to hold in every rotated frame, rotate the
third axis into the second. The same law then says \(A_2=B_2\); another
rotation gives \(A_1=B_1\). A frame-independent law was really the vector
equation
\begin{equation}
 A_i=B_i.
\end{equation}

Lorentz transformations extend the same argument to time. If \(A^0=B^0\) is
to hold in every inertial frame, boosts mix the time component with spatial
components and force the full equation \(A^\mu=B^\mu\). Good relativistic
laws must equate tensors of the same type.

The Newtonian mnemonic singled out \(T^{00}\), but the relativistic equation
must have the form
\begin{equation}
 G_{\mu\nu}=8\pi G\,T_{\mu\nu},
 \label{eq:field-schematic}
\end{equation}
where the symbol \(G_{\mu\nu}\) is initially only a placeholder. Its required
properties are now clear:
\begin{enumerate}
 \item it must be a symmetric rank-two tensor;
 \item it must be made from the metric, not from matter;
 \item it should contain two derivatives of the metric;
 \item its covariant divergence must vanish.
\end{enumerate}
The last condition follows from Equation~\eqref{eq:stress-conservation}.

\section{Searching Among Curvature Tensors}

The Christoffel symbols contain one derivative of the metric:
\begin{equation}
 \Gamma^\sigma_{\tau\nu}
 =
 \frac12 g^{\sigma\delta}
 \left(
 \partial_\tau g_{\delta\nu}
 +\partial_\nu g_{\delta\tau}
 -\partial_\delta g_{\tau\nu}
 \right).
 \label{eq:christoffel}
\end{equation}
They are not tensors, so they cannot by themselves provide the left-hand side.
The Riemann curvature tensor is a tensor:
\begin{equation}
 R^\rho{}_{\sigma\mu\nu}
 =
 \partial_\mu\Gamma^\rho_{\nu\sigma}
 -\partial_\nu\Gamma^\rho_{\mu\sigma}
 +\Gamma^\rho_{\mu\lambda}\Gamma^\lambda_{\nu\sigma}
 -\Gamma^\rho_{\nu\lambda}\Gamma^\lambda_{\mu\sigma}.
 \label{eq:riemann}
\end{equation}
Its first two terms contain second derivatives of the metric; its last two
contain products of first derivatives. This is exactly the differential
order suggested by the Newtonian equation.

The difficulty is index structure. \(R^\rho{}_{\sigma\mu\nu}\) has four
indices, while \(T_{\mu\nu}\) has two. Contracting a suitable upper and lower
index gives the Ricci tensor,
\begin{equation}
 R_{\mu\nu}=R^\lambda{}_{\mu\lambda\nu}.
 \label{eq:ricci}
\end{equation}
Other apparent contractions either vanish because of the symmetries of the
Riemann tensor or reproduce the same object up to a sign and index
convention. The Ricci tensor is symmetric, but it has less information than
the full Riemann tensor: it can vanish while the Riemann tensor remains
nonzero.

\begin{classroomqa}
\audiencequestion{Can one simply contract two lower indices of the Riemann
tensor? Does not a contraction require one covariant and one contravariant
index?}
\lecturerresponse{Exactly. An index must first be raised with the inverse
metric, or equivalently the contraction must include \(g^{\mu\nu}\).
Written as \(R^\lambda{}_{\mu\lambda\nu}\), the required upper-lower pair is
explicit.
\lecturetimestamp{01:00:09--01:00:33}}
\end{classroomqa}

Contracting the Ricci tensor again gives the curvature scalar,
\begin{equation}
 R=g^{\mu\nu}R_{\mu\nu}.
 \label{eq:ricci-scalar}
\end{equation}
Multiplying this scalar by \(g_{\mu\nu}\) returns another symmetric rank-two
tensor. We have therefore found two natural candidates:
\begin{equation}
 R_{\mu\nu},
 \qquad
 g_{\mu\nu}R.
\end{equation}

\begin{figure}[htbp]
\centering
\includegraphics[width=0.78\textwidth]{lecture_09_figure_03.png}
\caption{The board's two trial geometric tensors. At this stage both have
the correct rank and differential order, so each carries a question mark.}
\lectureframe{01:04:56}
\end{figure}

The exploratory alternatives are
\begin{align}
 R_{\mu\nu}&\stackrel{?}{=}8\pi G\,T_{\mu\nu},\\
 g_{\mu\nu}R&\stackrel{?}{=}8\pi G\,T_{\mu\nu}.
\end{align}
Neither is to be accepted merely because its indices fit. Conservation
provides the test.

\section{The Divergence Test}

If Equation~\eqref{eq:field-schematic} is correct, then
\begin{equation}
 D_\mu G^{\mu\nu}=0
\end{equation}
must follow identically, because \(D_\mu T^{\mu\nu}=0\). Begin with the
scalar-curvature candidate. The product rule gives
\begin{equation}
 D_\mu\!\left(g^{\mu\nu}R\right)
 =
 \left(D_\mu g^{\mu\nu}\right)R
 +g^{\mu\nu}D_\mu R.
 \label{eq:product-rule}
\end{equation}

\begin{figure}[htbp]
\centering
\includegraphics[width=0.84\textwidth]{lecture_09_figure_04.png}
\caption{Applying the covariant product rule to
\(D_\mu(g^{\mu\nu}R)\).}
\lectureframe{01:07:12}
\end{figure}

Metric compatibility says
\begin{equation}
 D_\mu g^{\alpha\beta}=0.
 \label{eq:metric-compatibility}
\end{equation}
This is an independent geometric property of the Levi-Civita connection, not
a consequence of energy-momentum conservation. Since \(R\) is a scalar,
\(D_\mu R=\partial_\mu R\). Equation~\eqref{eq:product-rule} therefore becomes
\begin{equation}
 D_\mu(g^{\mu\nu}R)
 =g^{\mu\nu}\partial_\mu R
 =D^\nu R,
\end{equation}
which is not generally zero. The candidate \(g^{\mu\nu}R\) fails.

The corresponding calculation for the Ricci tensor is somewhat longer. Its
result is the contracted Bianchi identity,
\begin{equation}
 D_\mu R^{\mu\nu}=\frac12 D^\nu R.
 \label{eq:contracted-bianchi}
\end{equation}
Thus \(R^{\mu\nu}\) also fails by itself. But the two failures have precisely
the ratio needed for cancellation:
\begin{equation}
 D_\mu\left(
 R^{\mu\nu}-\frac12 g^{\mu\nu}R
 \right)=0.
\end{equation}
This selects the Einstein tensor,
\begin{equation}
 \boxed{
 G_{\mu\nu}
 =
 R_{\mu\nu}-\frac12 g_{\mu\nu}R
 }.
 \label{eq:einstein-tensor}
\end{equation}

Within the assumptions being used here -- a local symmetric rank-two tensor
made from the metric with no more than two derivatives -- this combination is
fixed up to an overall coefficient and a possible metric term. The latter is
the cosmological-constant possibility, which returns near the end of the
lecture. The coefficient multiplying Equation~\eqref{eq:einstein-tensor} is
fixed by matching to Newtonian gravity.

\begin{classroomqa}
\audiencequestion{Is there a more directly physical way to obtain the field
equation, rather than trying candidates and rejecting them with the
continuity equation?}
\lecturerresponse{Local conservation is already a physical principle: energy
and momentum cannot disappear from a region without crossing its boundary.
There is also a more condensed derivation from an action principle for the
gravitational field. Varying that action produces the same equation, although
the calculation itself is not short.
\lecturetimestamp{01:17:14--01:18:41}}
\end{classroomqa}

\begin{classroomqa}
\audiencequestion{Could some other geometrical rank-two tensor satisfy the
same conservation law? Is the answer physical or mathematical, and would
Einstein have known it?}
\lecturerresponse{Under the assumptions of this construction, the
classification is mathematical. The available two-derivative curvature
objects descend from the Riemann tensor; the two rank-two candidates just
tested exhaust the relevant contractions. Finding their divergence-free
combination still requires a calculation. The principles are simple, but
managing the indices and signs is the labor.
\lecturetimestamp{01:18:42--01:21:44}}
\end{classroomqa}

\section{The Field Equation and Its Newtonian Limit}

The field equation is now
\begin{equation}
 \boxed{
 R_{\mu\nu}-\frac12 g_{\mu\nu}R
 =
 8\pi G\,T_{\mu\nu}
 }.
 \label{eq:einstein-equation}
\end{equation}
The left-hand side is geometric; the right-hand side describes matter,
radiation, and their flows. Their covariant divergences vanish together.

The factor \(8\pi G\) is not determined by the divergence argument. It is
determined by taking a weak, slowly varying gravitational field sourced by
slowly moving matter and requiring the time-time component to reproduce
Poisson's equation. In that regime \(T_{00}\simeq\rho\), while the other
components are suppressed by powers of \(v/c\). The geometric side reduces
to the second-derivative structure anticipated in
Equation~\eqref{eq:weak-mnemonic}. The match gives precisely the normalization
in Equation~\eqref{eq:einstein-equation}.

This limit also explains why Newtonian gravity appears to be sourced only by
mass density. In the full theory, every component of \(T_{\mu\nu}\) gravitates:
energy density, energy flow, momentum density, pressure, and shear. For
ordinary nonrelativistic matter, rest energy dominates. In relativistic
matter, radiation, or systems with large pressure, the remaining components
cannot be ignored.

The historical route was not a single effortless guess. The Ricci tensor
looks like the obvious first attempt, and the extra
\(-\tfrac12g_{\mu\nu}R\) is what restores compatibility with local
conservation. Once the tensor structure is fixed, the Newtonian limit
supplies the normalization.

\section{What Contraction Leaves Behind}

\begin{classroomqa}
\audiencequestion{When four indices of the Riemann tensor are contracted to
two, does the discarded information have physical content?}
\lecturerresponse{Yes. The Ricci tensor does not contain all of spacetime
curvature. Source-free gravitational fields provide the clearest examples:
the Ricci tensor can vanish while tidal curvature and gravitational waves
remain. The contraction keeps the part tied directly to the local source
equation, not every degree of freedom of the gravitational field.
\lecturetimestamp{01:21:57--01:23:10}}
\end{classroomqa}

The analogy with electromagnetism is useful. Maxwell's equations admit waves
in a region containing no charge or current. Likewise, Einstein's equation
admits nontrivial geometry in a region where \(T_{\mu\nu}=0\). We now derive
the simpler equation that holds there.

\section{Vacuum Means Ricci-Flat, Not Flat}

Set
\begin{equation}
 T_{\mu\nu}=0.
\end{equation}
Then Equation~\eqref{eq:einstein-equation} gives
\begin{equation}
 R_{\mu\nu}-\frac12g_{\mu\nu}R=0.
 \label{eq:vacuum-einstein}
\end{equation}
Contract with \(g^{\mu\nu}\). Since
\begin{equation}
 g^{\mu\nu}g_{\mu\nu}
 =\delta^\mu{}_\mu=4
\end{equation}
in four spacetime dimensions, the trace is
\begin{equation}
 R-\frac12(4)R=-R=0.
\end{equation}
Therefore
\begin{equation}
 R=0,
 \qquad
 R_{\mu\nu}=0.
 \label{eq:ricci-flat}
\end{equation}
The vacuum field equation is Ricci-flatness.

It is essential not to replace the words \emph{Ricci-flat} by \emph{flat}.
The Schwarzschild exterior contains no matter away from its central
singularity, and it satisfies \(R_{\mu\nu}=0\). Yet its Riemann tensor is
nonzero, and tidal forces are real. One can verify Ricci-flatness by
calculating all the Christoffel symbols and curvature components by hand,
though the calculation is a nuisance, or by letting a symbolic algebra system
do the bookkeeping. At the singularity the classical expressions diverge;
everywhere away from it, the vacuum equation holds.

Gravitational waves give a second example. They can propagate through a
source-free region while carrying genuine spacetime distortion. Thus
\begin{equation}
 R_{\mu\nu}=0
 \quad\not\Rightarrow\quad
 R^\rho{}_{\sigma\mu\nu}=0.
\end{equation}

\section{Keeping the Source: The Trace Equation}

The lecture nearly ends at the vacuum result, but an audience member asks to
repeat the contraction without discarding the source.

\begin{classroomqa}
\audiencequestion{What follows if we perform the same contraction while
leaving \(T_{\mu\nu}\) in the equation?}
\lecturerresponse{The trace relates the scalar curvature directly to the
trace of the energy-momentum tensor, and substituting that relation back gives
an equivalent Ricci-tensor form of the field equation.
\lecturetimestamp{01:29:28--01:32:20}}
\end{classroomqa}

Define
\begin{equation}
 T=g^{\mu\nu}T_{\mu\nu}.
\end{equation}
Taking the trace of Equation~\eqref{eq:einstein-equation} gives
\begin{align}
 R-\frac12(4)R&=8\pi G\,T,\\
 -R&=8\pi G\,T,
\end{align}
so
\begin{equation}
 \boxed{R=-8\pi G\,T}.
 \label{eq:trace-relation}
\end{equation}
Substitution back into the field equation yields
\begin{equation}
 \boxed{
 R_{\mu\nu}
 =
 8\pi G
 \left(
 T_{\mu\nu}-\frac12g_{\mu\nu}T
 \right)
 }.
 \label{eq:ricci-form}
\end{equation}

\begin{figure}[htbp]
\centering
\includegraphics[width=0.82\textwidth]{lecture_09_figure_05.png}
\caption{The trace relation and the equivalent Ricci-tensor form of
Einstein's equation.}
\lectureframe{01:32:25}
\end{figure}

\begin{classroomqa}
\audiencequestion{Does the scalar \(T\) have a physical significance?}
\lecturerresponse{It is the trace of the energy-momentum tensor. For
electromagnetic radiation it vanishes, and the same is true for an ideal gas
of massless particles such as photons. Massive matter generally has a
nonzero trace. Through Equation~\eqref{eq:trace-relation}, that trace fixes
the scalar curvature.
\lecturetimestamp{01:32:20--01:33:12}}
\end{classroomqa}

\section{How Much Curvature Does Ricci Remember?}

\begin{classroomqa}
\audiencequestion{Is there a simple physical or geometric picture of the
Ricci tensor and Ricci scalar, comparable to parallel transport around a
small loop for the Riemann tensor?}
\lecturerresponse{No comparably transparent picture is offered here. The
Riemann tensor directly records the rotation produced by transport around a
small loop. Its contractions are simpler algebraically, but their individual
geometric meanings are harder to visualize. It is better not to invent a
misleading picture merely because one would be convenient.
\lecturetimestamp{01:33:14--01:34:13}}
\end{classroomqa}

The amount of information lost in the contraction depends on dimension. In
four dimensions, the Riemann tensor contains curvature information not
present in the Ricci tensor. In three dimensions, the Riemann tensor is
determined by the Ricci tensor. In two dimensions, the scalar curvature
contains all the independent curvature information. Four-dimensional
gravity therefore has room for source-free curvature degrees of freedom that
are invisible to \(R_{\mu\nu}\).

\begin{classroomqa}
\audiencequestion{If different Riemann tensors can have the same Ricci tensor,
how does that missing information appear physically?}
\lecturerresponse{A given source does not by itself specify every possible
gravitational configuration. Boundary and initial data can include
gravitational radiation. In the weak-field approximation one may think of
starting with a sourced solution and adding a free gravitational wave. The
full nonlinear theory does not permit arbitrary superposition, but the
example correctly identifies the missing information: propagating curvature
not fixed locally by \(T_{\mu\nu}\).
\lecturetimestamp{01:37:05--01:38:59}}
\end{classroomqa}

The vacuum examples now fit together. Schwarzschild curvature records the
tidal field outside a mass; gravitational waves record propagating free
curvature. Both can satisfy \(R_{\mu\nu}=0\) without being flat.

\section{The Cosmological Constant}

\begin{classroomqa}
\audiencequestion{Is there no additional fudge factor, such as a
cosmological constant?}
\lecturerresponse{A cosmological constant is allowed. It can be written on
the geometric side as a constant times \(g_{\mu\nu}\), or shifted to the
right and interpreted as a vacuum contribution to the energy-momentum
tensor. The two descriptions are the same equation with different
bookkeeping.
\lecturetimestamp{01:39:01--01:40:44}}
\end{classroomqa}

With a cosmological constant, one common convention is
\begin{equation}
 G_{\mu\nu}+\Lambda g_{\mu\nu}
 =
 8\pi G\,T_{\mu\nu}.
\end{equation}
Moving the extra term to the right defines an effective vacuum stress tensor,
\begin{equation}
 T_{\mu\nu}^{(\Lambda)}
 =
 -\frac{\Lambda}{8\pi G}g_{\mu\nu}.
\end{equation}
Because metric compatibility gives \(D_\mu g^{\mu\nu}=0\), the extra term is
automatically compatible with the conservation law. This is the metric term
left open by the earlier uniqueness statement.

\section{Mercury and the Weak-Field Calculation}

\begin{classroomqa}
\audiencequestion{Did Einstein use this equation to calculate Mercury's
orbit, and where can one see how that calculation was done?}
\lecturerresponse{Einstein did not yet have Schwarzschild's exact solution,
but he had a sufficiently accurate large-distance approximation to the
exterior field of the Sun. The Sun's radius and Mercury's orbit are enormous
compared with the Sun's Schwarzschild radius, so relativistic corrections are
small. He could begin with a Keplerian orbit and compute the correction
perturbatively. Schwarzschild obtained the exact exterior solution soon
afterward, making the same calculation systematic.
\lecturetimestamp{01:40:46--01:43:48}}
\end{classroomqa}

Perturbation theory means exactly this: start with a problem already solved,
then calculate the effect of a small correction. Newton's inverse-square
force produces closed Keplerian ellipses. General relativity adds a small
departure from the Newtonian radial equation, so the ellipse does not close
after one revolution; its perihelion advances.

The lecture closes with the curious specialness of closed central-force
orbits. The inverse-square force is exceptional in making the Keplerian
orbits close without precession. More generally, Bertrand's theorem also
identifies the isotropic harmonic oscillator as the other central potential
for which every bound orbit is closed. The small relativistic correction
removes that Newtonian degeneracy, turning Mercury's precession into an
observable signature of the field equation derived in this chapter.
